


The developments in \cref{sec:timo} are most recent, and the results were much more general than I expected. Consequently \cref{sec:} and \cref{sec:gensec} need to be extensively revised/updated.

For \cref{sec:timo}, first read the summary \cref{sec:trace-summary}.

Primary outcomes of \cref{sec:timo} are:
\begin{itemize}
    \item In a fiber section we have no way to use the classical shear correction factors, and instead we have to produce the same effect with the matrices I currently denote by \(\mathbf{L}_{\cdot}\).
    \item \cref{sec:timo} shows that when these are derived by requiring fiber strains be exact under Saint-Venant loads, one finds 18 shear-correction-like parameters. 
    \item There are several ways to formulate these 18 constants which all produce continuum-exact fiber strains. In the current organization, \cref{sec:timo} only develops these general results and doesn't derive any of these formulations in particular. It just shows derives what is necessary to implement a section that supports these parameters.
    \item Two particular formulations are derived in \cref{sec:trace-cowper} and \cref{sec:trace-energy}.
    \item The meaning of these parameters relates to the 3D nature of the boundary conditions. 
\end{itemize}